\chapter{Zadatak 1e --- Lagrangeov dualni problem}
\label{ch:z1-dual}

\section{Što se ovdje traži}

Izvesti dualni problem \emph{analitički}, riješiti ga numerički i usporediti
dualne varijable s onima iz (b) i (d).

\begin{zadatak}{Zadatak 1, podzadatak e)}
e) Postavite i riješite Lagrangeov dualan problem (koristite MATLAB za numeričko
rješavanje dualnog problema). Usporedite vrijednost izračunatih dualnih varijabli
s onima dobivenim pod (d) i (b).
\end{zadatak}

\begin{predznanje}
\begin{itemize}[nosep,leftmargin=*]
  \item Lagrangeova funkcija \eqref{eq:z1b-lagrangian} i množitelji $\mu \geq 0$.
  \item $Q^{-1}$ postoji jer je $Q \pd 0$ --- \S\ref{sec:inverz}, \eqref{eq:alat-Q-inv}.
  \item Jaka dualnost i Slaterov uvjet --- \S\ref{sec:z1a-slater}.
  \item Značenje $\inf$ nasuprot $\min$ --- \S\ref{sec:oznake-optimizacije}.
\end{itemize}
\end{predznanje}

\section{Definicije s predavanja}
\label{sec:z1e-definicije}

\begin{izvor}{Predavanje 4 --- Lagrangeova dualnost, slajd 6/26 (PDF str.~10)}
\slajd[0.86]{p04-lagrangian-i-dual}

\textbf{Transkripcija.} \emph{Terminologija i neka svojstva.}
\begin{itemize}[nosep,leftmargin=*]
  \item Lagrangian (Lagrangeova funkcija):
        $\Lag(x,\lambda,\mu) \coloneqq f(x) + \sum_{i=1}^{m}\lambda_i h_i(x) + \sum_{j=1}^{p}\mu_j g_j(x)$
  \item Lagrangeova dualna funkcija cilja:
        $\dual(\lambda,\mu) \coloneqq \inf_x \Lag(x,\lambda,\mu)$
  \item Lagrangeov dualan problem:
        $\dopt = \sup_{\{\lambda,\mu \mid \mu \geq 0\}} \dual(\lambda,\mu)$
  \item Primarni problem: $\popt = \inf_{\{x \mid g(x) \leq 0,\, h(x) = 0\}} f(x)$
  \item $\lambda_i$ je Lagrangeov množitelj (dualna varijabla) pridružen
        ograničenju $h_i(x) = 0$;
        $\mu_j$ je Lagrangeov množitelj pridružen ograničenju $g_j(x) \leq 0$.
\end{itemize}

\textbf{Tri stvari koje treba uočiti u ovoj definiciji:}
\begin{enumerate}[nosep,leftmargin=*]
  \item Dualna funkcija nastaje \emph{minimizacijom po $x$ bez ikakvih
        ograničenja} --- ograničenja su ušla u Lagrangian i time nestala iz skupa
        po kojem se traži infimum.
  \item Dualna funkcija je \emph{uvijek konkavna}, bez obzira na to je li primarni
        problem konveksan: infimum familije afinih funkcija po $(\lambda,\mu)$ je
        konkavna funkcija. Zato je dualni problem \emph{uvijek} konveksan.
  \item Jedina preostala ograničenja u dualnom problemu su $\mu \geq 0$.
\end{enumerate}
\end{izvor}

\section{Izvod dualne funkcije za QP}
\label{sec:z1e-izvod}

\begin{izvor}{Vježbe 3 --- Primjer 2: Dualnost, KKT, osjetljivost, slajd 30/36 (PDF str.~37)}
\slajd[0.86]{v3-primjer2-dual-izvod}

\textbf{Transkripcija.} \emph{d) Postavi i riješi dualan problem.} Primaran
problem:
\[
\min_x \ \tfrac12 x\trans Qx + c\trans x \quad \text{uz ogr.} \quad Gx - h \leq 0 .
\]
Uočimo: $Q \pd 0$, pa stoga i $Q^{-1}$ postoji ($Q$ nema svojstvenu vrijednost
jednaku nuli jer su sve strogo veće od nule).
\[
\Lag(x,\mu) = \tfrac12 x\trans Qx + c\trans x + \mu\trans(Gx - h)
            = \tfrac12 x\trans Qx + (c + G\trans\mu)\trans x - \mu\trans h
\]
Dualan problem:
\[
\max_{\mu \geq 0} \ -\tfrac12 \mu\trans G Q^{-1} G\trans \mu
+ \bigl(-c\trans Q^{-1}G\trans - h\trans\bigr)\mu - \tfrac12 c\trans Q^{-1} c .
\]

\textbf{Ovo je opći rezultat za QP} --- profesor ga izvodi u simbolima, bez
brojeva, pa vrijedi i za naš zadatak. Sve što nam preostaje jest uvrstiti naše
$Q$, $c$, $G$, $h$.
\end{izvor}

\begin{korak}[Izvod, korak po korak]
\begin{enumerate}[leftmargin=*]
  \item \textbf{Grupiraj po $x$.} Uvrštavanjem $g(x) = Gx - h$ u definiciju
        Lagrangiana i skupljanjem linearnih članova:
        \[
          \Lag(x,\mu) = \tfrac12 x\trans Q x + \underbrace{(c + G\trans\mu)}_{\eqqcolon\, w}{}\trans x - \mu\trans h .
        \]
        Uveli smo kraticu $w \coloneqq c + G\trans\mu$; ona ovisi o $\mu$, ali ne o $x$.
  \item \textbf{Infimum po $x$ se \emph{postiže}.} Kako je $Q \pd 0$, funkcija
        $x \mapsto \Lag(x,\mu)$ je strogo konveksna po $x$, pa infimum nije samo
        donja granica nego se doseže u jedinstvenoj točki. Ta je točka rješenje
        uvjeta $\grad_x \Lag = 0$:
        \[
          Q\tilde{x} + w = 0 \;\Longrightarrow\; \tilde{x} = -Q^{-1}w .
        \]
  \item \textbf{Uvrsti $\tilde{x}$ natrag.} Uz $\tilde{x} = -Q^{-1}w$ i
        simetričnost $Q^{-1}$:
        \begin{align*}
          \tfrac12 \tilde{x}\trans Q \tilde{x}
            &= \tfrac12 w\trans Q^{-1} Q Q^{-1} w = \tfrac12 w\trans Q^{-1} w, \\
          w\trans \tilde{x} &= -w\trans Q^{-1} w,
        \end{align*}
        pa je
        \[
          \dual(\mu) = \tfrac12 w\trans Q^{-1}w - w\trans Q^{-1}w - \mu\trans h
                     = -\tfrac12 w\trans Q^{-1}w - \mu\trans h .
        \]
  \item \textbf{Raspiši $w = c + G\trans\mu$.}
        \[
          w\trans Q^{-1} w = c\trans Q^{-1}c + 2\,c\trans Q^{-1}G\trans\mu + \mu\trans G Q^{-1}G\trans \mu ,
        \]
        odakle
        \[
          \dual(\mu) = -\tfrac12 \mu\trans G Q^{-1}G\trans\mu
                       + \bigl(-c\trans Q^{-1}G\trans - h\trans\bigr)\mu
                       - \tfrac12 c\trans Q^{-1}c .
        \]
\end{enumerate}
\end{korak}

\begin{formula}{Dualni problem}
\begin{equation}
\max_{\mu \geq 0}\ \dual(\mu),
\qquad
\dual(\mu) = -\tfrac{1}{2}\mu\trans \underbrace{G Q^{-1} G\trans}_{\text{PSD}}\mu
+ \bigl(-c\trans Q^{-1}G\trans - h\trans\bigr)\mu
- \tfrac{1}{2}c\trans Q^{-1}c .
\label{eq:z1e-dual}
\end{equation}
Matrica $-GQ^{-1}G\trans$ je negativno semidefinitna, pa je $\dual$ konkavna
kvadratna funkcija. Maksimizacija konkavne funkcije uz ograničenje $\mu \geq 0$
je konveksan problem --- opet kvadratni program, samo u varijabli $\mu$.
\end{formula}

\section{Dualni problem u brojevima}
\label{sec:z1e-brojevi}

Uvrštavanjem $Q^{-1} = \tfrac{100}{791}\begin{bsmallmatrix}4 & 0{,}3\\ 0{,}3 & 2\end{bsmallmatrix}$,
$G$, $c$ i $h$ dobiva se, egzaktno u razlomcima:
\begin{equation}
G Q^{-1} G\trans = \frac{1}{791}\begin{bmatrix} 660 & -200 \\ -200 & 540 \end{bmatrix}
\approx \begin{bmatrix} 0{,}834387 & -0{,}252845 \\ -0{,}252845 & 0{,}682680 \end{bmatrix},
\end{equation}
\begin{equation}
-c\trans Q^{-1}G\trans - h\trans = \frac{1}{791}\begin{bmatrix} 8080 & -3623 \end{bmatrix}
\approx \begin{bmatrix} 10{,}214918 & -4{,}580278 \end{bmatrix},
\qquad
\tfrac12 c\trans Q^{-1}c = \frac{1520}{791} .
\end{equation}

\begin{korak}[Uočljiva veza: linearni član dualne funkcije]
Linearni koeficijent $-c\trans Q^{-1}G\trans - h\trans$ nije proizvoljan broj.
Vrijedi
\[
-G Q^{-1}c - h = G x^{0} - h = g(x^{0}),
\]
jer je $x^{0} = -Q^{-1}c$ neograničeni minimum iz Slučaja 1. Doista, prva
komponenta je $8080/791 = 10{,}214918$ --- upravo vrijednost $g_1(x^{0})$ zbog
koje je Slučaj 1 odbačen (\S\ref{sec:z1b-slucaj1}). Dualna funkcija u nuli ima
nagib jednak vrijednostima ograničenja u neograničenom minimumu.
\end{korak}

\begin{korak}[Provjera: $\dual(\muopt) = \popt$]
Uz $\muopt = (404/33,\; 0)$ druga komponenta ne doprinosi, pa je
\[
\dual(\muopt) = -\frac12 \cdot \frac{660}{791}\left(\frac{404}{33}\right)^{2}
+ \frac{8080}{791}\cdot\frac{404}{33} - \frac{1520}{791}
= \frac{2000}{33} = 60{,}606061 .
\]
To je točno $\popt$ iz \eqref{eq:z1b-rjesenje}: dualni procjep je \textbf{egzaktno
nula}, kako i mora biti kad vrijedi jaka dualnost.
\end{korak}

\subsection{Zašto je maksimum na rubu $\mu_2 = 0$}

\begin{korak}
Da nema ograničenja $\mu \geq 0$, maksimum konkavne kvadratne funkcije
\eqref{eq:z1e-dual} bio bi u stacionarnoj točki
$\bar{\mu} = (GQ^{-1}G\trans)^{-1}\bigl(-GQ^{-1}c - h\bigr)$. Izračunato:
\begin{equation}
\bar{\mu} = \begin{bmatrix} 23/2 \\ -49/20 \end{bmatrix}
          = \begin{bmatrix} 11{,}50 \\ -2{,}45 \end{bmatrix}.
\end{equation}
To su \emph{točno} množitelji koje smo dobili u Slučaju 4
(\S\ref{sec:z1b-slucaj4}). Druga komponenta je negativna, dakle
$\bar{\mu}$ nije dopustiv za dualni problem, pa se maksimum postiže na rubu
$\mu_2 = 0$. Ondje je
\begin{equation}
\grad \dual(\muopt) = \begin{bmatrix} 0 \\ -49/33 \end{bmatrix}
                     = \begin{bmatrix} g_1(\xopt) \\ g_2(\xopt) \end{bmatrix},
\end{equation}
što je opće svojstvo: gradijent dualne funkcije jednak je vrijednostima
ograničenja u pripadnom $\tilde{x}(\mu)$. Prva komponenta je nula jer je $g_1$
aktivno; druga je negativna, pa bi povećanje $\mu_2$ smanjilo $\dual$ --- rub je
optimalan.
\end{korak}

\begin{pojam}{Ista činjenica, tri lica}
Slučaj 4 KKT sustava, unutarnji maksimum dualne funkcije i vrh dopustivog skupa
tri su opisa iste točke $(11{,}5,\; -2{,}45)$. U primarnom svijetu odbacuje se
zbog $\mu_2 < 0$; u dualnom, zato što je izvan dopustivog skupa $\mu \geq 0$.
Ako to na obrani povežeš, pokazao si da razumiješ dualnost, a ne da si zapamtio
dva odvojena računa.
\end{pojam}

\section{Kako je to riješeno u seminaru}
\label{sec:z1e-seminar}

\begin{seminar}{Seminar v1.2, §2.6 --- ,,(e) Lagrangeov dualni problem``}
Budući da je $Q \succ 0$, matrica $Q^{-1}$ postoji, a Lagrangeova funkcija
\[
\mathcal{L}(x, \mu) = \tfrac{1}{2} x^{\mathsf{T}} Q x + (c + G^{\mathsf{T}}\mu)^{\mathsf{T}} x - \mu^{\mathsf{T}} h
\]
strogo je konveksna po $x$. Infimum se postiže u točki u kojoj je
$\nabla_x \mathcal{L} = 0$:
\[
Q\tilde{x} + c + G^{\mathsf{T}} \mu = 0 \implies \tilde{x} = -Q^{-1}(c + G^{\mathsf{T}} \mu).
\]
Uvrštavanjem $\tilde{x}$ natrag u $\mathcal{L}$ dobiva se dualna funkcija
\[
\theta(\mu) = \inf_{x} \mathcal{L}(x, \mu)
= -\tfrac{1}{2} \mu^{\mathsf{T}} G Q^{-1} G^{\mathsf{T}} \mu
+ \left( -c^{\mathsf{T}} Q^{-1} G^{\mathsf{T}} - h^{\mathsf{T}} \right) \mu
- \tfrac{1}{2} c^{\mathsf{T}} Q^{-1} c,
\]
a dualni problem glasi $\max_{\mu \geq 0} \theta(\mu)$.

Matrica $-G Q^{-1} G^{\mathsf{T}}$ je negativno semidefinitna, pa je $\theta$
konkavna kvadratna funkcija, a dualni problem je maksimizacija konkavne funkcije
uz konveksna ograničenja, dakle konveksan optimizacijski problem. Riječ je opet o
kvadratnom programu, samo u varijabli $\mu$, pa se rješava istom funkcijom
\texttt{quadprog} primijenjenom na $-\theta(\mu)$ uz donju granicu $\mu \geq 0$.

\begin{lstlisting}[style=matlabstyle]
Qi = inv(Q);
Hd = G*Qi*G';  Hd = 0.5*(Hd + Hd');    % Hessian dualnog QP-a
fd = (c'*Qi*G' + h')';
kd = 0.5*(c'*Qi*c);                    % konstantni clan
[mu_d, negd, ef_d] = quadprog(Hd, fd, [], [], [], [], [0;0], [], [], opts);
d_opt = -(negd + kd);
\end{lstlisting}

\begin{lstlisting}[style=outputstyle]
exitflag        = 1
mu*             = [12.242424; 0.000000]
d*  = theta(mu*)= 60.606061
dualni procjep p* - d* = 2.480e-10
\end{lstlisting}

Dualne varijable $\mu^*$ identične su Lagrangeovim množiteljima dobivenim
analitički pod (b) i onima koje je \texttt{quadprog} vratio pod (d). Kako je
Slaterov uvjet zadovoljen, vrijedi jaka dualnost, što potvrđuje i numerički
dualni procjep $p^* - d^* = 2.48 \cdot 10^{-10}$. Maksimum konkavne funkcije
$\theta$ postiže se na rubu $\mu_2 = 0$, što odgovara neaktivnosti drugog
primarnog ograničenja.
\end{seminar}

\subsection{Zašto kod izgleda baš tako}

\begin{longtable}{@{}p{0.34\linewidth} p{0.62\linewidth}@{}}
\toprule
\textbf{Redak} & \textbf{Razlog} \\
\midrule
\endhead
\texttt{Hd = 0.5*(Hd + Hd')} & Simetrizacija. Matematički je $GQ^{-1}G\trans$
simetrična, ali u aritmetici s pomičnim zarezom razlikuje se od svog transponata
u zadnjim bitovima, a \texttt{quadprog} traži simetričan ulaz. \\[5pt]
\texttt{fd = (c'*Qi*G' + h')'} & Predaje se \emph{minus} linearni član iz
\eqref{eq:z1e-dual}, jer se minimizira $-\dual$. Predznak je već ugrađen. \\[5pt]
\texttt{[0;0]} kao \texttt{lb} & Jedino ograničenje dualnog problema, $\mu \geq 0$,
zadaje se kao donja granica, ne kao redak matrice $A$. \\[5pt]
\texttt{d\_opt = -(negd + kd)} & \texttt{quadprog} vraća minimum od $-\dual$ bez
konstantnog člana, pa se konstanta vraća i predznak okreće. \\
\bottomrule
\end{longtable}

\begin{zamka}[Zamka: prešućeni konstantni član]
Konstanta $-\tfrac12 c\trans Q^{-1}c = -1520/791 \approx -1{,}9216$ ne utječe na
to \emph{gdje} je maksimum, pa je rješavaču nije potrebno predati. Ali utječe na
\emph{vrijednost} $\dopt$. Ako se zaboravi vratiti, dualni procjep ispadne
$\approx 1{,}92$ umjesto nule i cijeli zaključak o jakoj dualnosti pada.
\end{zamka}

\section{Modelski odgovori}
\label{sec:z1e-odgovori}

\begin{pitanje}{Zašto je dualni problem uvijek konveksan, čak i kad primarni nije?}
\textbf{Odgovor.} Jer je dualna funkcija infimum familije funkcija koje su afine
u $(\lambda,\mu)$: za svaki fiksni $x$, preslikavanje
$\mu \mapsto f(x) + \mu\trans g(x)$ je afino. Infimum bilo koje familije
konkavnih (ovdje afinih) funkcija je konkavan, a maksimizacija konkavne funkcije
uz $\mu \geq 0$ je konveksan problem. Konveksnost dualnog problema ne ovisi ni o
čemu u primarnome.
\end{pitanje}

\begin{pitanje}{Koliko bi iznosio dualni procjep da problem nije konveksan?}
\textbf{Odgovor.} Bio bi nenegativan, ali ne nužno nula --- vrijedila bi samo
slaba dualnost $\dopt \leq \popt$. Ovdje je nula jer je problem konveksan i
Slaterov uvjet zadovoljen. Numerički se dobiva $2{,}48 \cdot 10^{-10}$, što je
tolerancija rješavača, a ne stvaran procjep.
\end{pitanje}

\begin{pitanje}{Čemu uopće služi dualni problem ako primarni već znamo riješiti?}
\textbf{Odgovor.} Trima stvarima. Prvo, daje \emph{certifikat}: svaki dopustivi
$\mu \geq 0$ daje donju granicu na $\popt$, pa se optimalnost može dokazati bez
rješavanja primarnog problema. Drugo, dualne varijable imaju izravno inženjersko
značenje --- osjetljivost, o čemu je podzadatak (f). Treće, za velike probleme
dual može biti mnogo manji od primarnog, jer ima onoliko varijabli koliko primar
ima ograničenja.
\end{pitanje}

\begin{izgovor}
,,Lagranžijan sam grupirao po iks, pa je linearni član ce plus ge transponirano
mi. Budući da je ku pozitivno definitna, Lagranžijan je strogo konveksan po iks,
pa se infimum stvarno postiže, i to u točki iks tilda jednako minus ku inverz
puta ce plus ge transponirano mi.

Uvrštavanjem natrag dobiva se dualna funkcija: minus jedna polovina mi
transponirano ge ku inverz ge transponirano mi, plus linearni član, minus
konstanta jedna polovina ce transponirano ku inverz ce. Matrica ispred kvadratnog
člana je negativno semidefinitna, pa je dualna funkcija konkavna, i dualni
problem je maksimizacija konkavne funkcije uz mi veće ili jednako nuli --- dakle
opet kvadratni program.

Riješio sam ga u MATLAB-u i dobio mi jednako dvanaest zarez dva četiri i nula,
identično množiteljima iz podzadataka b) i d), s optimalnom vrijednošću šezdeset
zarez šest nula šest. Dualni procjep je dva puta deset na minus deset, dakle
numerička nula --- u skladu s jakom dualnošću.

Zanimljivo je da bi bezuvjetni maksimum dualne funkcije bio u jedanaest i pol,
minus dva zarez četiri pet. To su točno množitelji iz slučaja u kojem su oba
ograničenja aktivna, a odbacili smo ga upravo zato što je drugi množitelj
negativan. Ista činjenica, gledana iz primarnog i iz dualnog problema.``
\end{izgovor}
